\subsection*{作者标准术语：原书定义 1.5.2、1.5.4 与注记 1.5.5}
\noindent\textit{整理说明：以下原文从同一版本 PDF 物理页 45--46（书内页 37--38）转录；锥范畴采用通常定义。}
\par\medskip
\def\CorpusNumber{1.5.2}\begin{definition}\label{def:create-limit}
给定 $F: \mathcal{C}\to\mathcal{D}$ 和函子 $\alpha:I\to\mathcal{C}$, 考虑锥范畴之间的函子 $(\alpha/\Delta)\to(F\alpha/\Delta)$.
\begin{itemize}
\item 称 $F$ \textbf{保}此 $\varinjlim$, 如果 $(\alpha/\Delta)\to(F\alpha/\Delta)$ 映始对象 (假如存在) 为始对象.
\item 称 $F$ \textbf{返}此 $\varinjlim$, 如果仅 $(\alpha/\Delta)$ 的始对象方能被映为 $(F\alpha/\Delta)$ 的始对象.
\item 称 $F$ \textbf{生}此 $\varinjlim$, 如果一旦 $\varinjlim F\alpha$ 存在, 则 $\varinjlim\alpha$ 存在, 而且此时 $F$ 保此 $\varinjlim$ 也返此 $\varinjlim$.
\end{itemize}
对于 $\varprojlim$ 同样有相应的概念, 以 $\mathcal{C}^{\opp}$ 和 $\mathcal{D}^{\opp}$ 代 $\mathcal{C}$ 和 $\mathcal{D}$ 即可相互过渡.
\end{definition}
\def\CorpusNumber{1.5.4}\begin{definition}\label{def:conservative}
满足以下条件的函子 $F$ 称为\textbf{保守}的: 态射 $f\in\Mor(\mathcal{C})$ 为同构当且仅当 $Ff\in\Mor(\mathcal{D})$ 为同构.
\end{definition}
\def\CorpusNumber{1.5.5}\begin{remark}\label{rem:conservative-reflect}
若 $F$ 是保守函子, 则只要 $\varinjlim\alpha$ 存在, 而且 $F$ 保此 $\varinjlim$, 则 $F$ 也必然返此 $\varinjlim$. 缘由如下: 设 $L\in\Obj((\alpha/\Delta))$ 被映为 $(F\alpha/\Delta)$ 的始对象; 考虑 $(\alpha/\Delta)$ 中的唯一态射 $\varinjlim\alpha\to L$, 由于 $F$ 保此 $\varinjlim$, 它在 $F$ 之下的像必为同构, 从而保守条件确保 $\varinjlim\alpha\to L$ 也是同构.

在此前提下, 关于 $F$ 生 $\varinjlim$ 的定义中可以去除返 $\varinjlim$ 的条件.
\end{remark}
